\documentclass[aspectratio=169]{ctexbeamer}
\usepackage[T1]{fontenc}
\usepackage{latexsym,amsmath,xcolor,multicol,booktabs,calligra}
\usepackage{graphicx,listings,stackengine}
\usepackage{hyperref}

\usepackage{PekingU}

% 去掉每节/子节前面的自动目录页
\AtBeginSection[]{}
\AtBeginSubsection[]{}

% ===== 每节课改这里 =====
\author{Cap1taL}
\title{模拟赛1讲解}
\subtitle{meow}
\date{2026年8月}
% ========================

% 代码高亮设置
\definecolor{deepblue}{rgb}{0,0,0.5}
\definecolor{deepred}{rgb}{0.6,0,0}
\definecolor{deepgreen}{rgb}{0,0.5,0}
\definecolor{halfgray}{gray}{0.55}

\lstset{
    basicstyle=\ttfamily\small,
    keywordstyle=\bfseries\color{deepblue},
    emphstyle=\ttfamily\color{deepred},
    stringstyle=\color{deepgreen},
    numbers=left,
    numberstyle=\small\color{halfgray},
    rulesepcolor=\color{red!20!green!20!blue!20},
    frame=shadowbox,
}

\begin{document}

\kaishu

% 封面
\begin{frame}
    \titlepage
\end{frame}

% 目录
% \begin{frame}{目录}
%     \tableofcontents[sectionstyle=show,subsectionstyle=show/shaded/hide,subsubsectionstyle=show/shaded/hide]
% \end{frame}

% ===== 从这里开始写你的内容 =====

\section{Anon 的花园}

\subsection{Anon 的花园}
\begin{frame}{题意}
    对于所有位置，找到尽可能大的正方形，使得内部颜色种类数不超过 $k$

    $n\le 500,\; k,a_{i,j}\le n$
\end{frame}

\begin{frame}{Solution}
    \begin{itemize}
        \item 考虑暴力做法，直接枚举所有正方形，然后遍历正方形维护颜色数量，复杂度为 $O(n^5)$，固定端点后精细实现可以做到 $O(n^4)$ ，可以通过 $n\le 50$，此处不再赘述
        \item 考虑 $a\le 2$，只需要进一步讨论 $k=1$ 和 $k\ge 2$，前者需要找同色最大正方形，后者则总是最大的正方形
        \item 考虑 $k=2$，会发现等价于我们固定住端点扩展长度，当 “出现第三个颜色” 时非法
        \item 用各类方法实现 $k=2$ 后，会自然引导我们尝试刻画 “每个颜色什么时候出现”
    \end{itemize}
\end{frame}

\begin{frame}{Solution}
    \begin{itemize}
        \item 考虑正解，观察到值域并不大，我们可以接受与值域有关的复杂度
        \item 先对每个位置、每种颜色处理 $f_{i,j,c}$ 表示以 $(i,j)$ 为左上角，边长多长时第一次出现 $c$ 这个颜色。可以从 $(i+1,j+1),(i,j+1),(i+1,j)$ 转移到 $(i,j)$。
        \item 有了 $f$ 之后，我们可以进一步求出 $g_{i,j,L}$ 表示有多少个 $c$ 满足 $f_{i,j,c}=L$。可以直接用 $f$ 计算，这一部分是简单的。
        \item 基于此，我们就可以枚举每个左上角端点 $(i,j)$，然后线性扫描 $g_{i,j,L}$，使得前缀和不超过 $k$ 的最大 $L$ 即为答案。复杂度 $O(n^3)$
        \item 具体实现上，开两个 $O(n^3)$ 的数组可能超过内存限制。可以考虑将 int 更换为 short，或者更好的做法是不完整存下 $f_{i,j,c}$，而是对每个颜色处理出 $f_{i,j}$ 后直接计算 $g$ 数组。
    \end{itemize}
\end{frame}

\section{Soyo 的定向}
\subsection{Soyo 的定向}
\begin{frame}{题面}
    一棵树，有边权，边权为一个排列。要求给无向边定向，使一个所有可能的路径都边权递减，问定向方案数

    $n\le 3\times 10^5$
\end{frame}

\begin{frame}{Solution}
    \begin{itemize}
        \item $n\le 15$ 的部分分可以直接暴力枚举定向方案，复杂度为指数级
        \item $n\le 3000$ 给一些可能的平方做法，或者 dp 实现不精细的做法
        \item 考虑链，容易发现决策满足很好的无后效性，可以直接做线性 DP，设 $f_{i,0/1}$ 表示 $i$ 左侧的边的朝向，讨论边权大小关心转移即可
        \item 考虑菊花，此时可能的路径长度只有 $ 2 $。由于要求路径权值递增，因此不能存在 “较大的边朝向中心，较小的边朝向外侧” 的情况。
        \item 更具体的，若我们把菊花图的边按照边权升序排序，那么定向方案一定是前缀朝向中心，后缀远离中心
    \end{itemize}
\end{frame}

\begin{frame}{Solution}
    \begin{itemize}
        \item 链和菊花已经给了我们足够多的提示
        \item 把我们刚才的算法上树，即设 $f_{u,0/1}$ 表示 $u$ 到父亲边为 $0/1$ 的情况下整个子树的方案数
        \item 转移时模仿菊花图，给每个点所有的边排序，之后用子节点更新父节点即可
        \item 具体实现时可以把连向父亲的边和与孩子的边一起排序，然后扫描记录 flag 来决定转移到 $f_{u,0}$ 还是 $f_{u,1}$
        \item 复杂度 $O(n\log n)$，瓶颈在给边排序
    \end{itemize}
\end{frame}


\section{Tomori 的石头}
\subsection{Tomori 的石头}
\begin{frame}{题面}
    一个序列，每次随机合并相邻两个数，合并结果是他们的平均数，直到序列长度为 $ 1 $。求最终结果的期望

    $n\le 10^6$
\end{frame}

\begin{frame}{Solution}
    \begin{itemize}
        \item $n$ 较小的部分分留给分类讨论或者暴力枚举的做法
        \item 观察题目合并的形式，实际上答案形如 $\sum \frac{a_i}{2^{k_i}}$，其中 $k_i$ 的数值与合并顺序有关
        \item 根据期望的线性性，我们有
        $$
            E\left(\sum \dfrac{a_i}{2^{k_i}}\right)=\sum E\left(\dfrac{a_i}{2^{k_i}}\right)=\sum a_i\times E\left(\dfrac{1}{2^{k_i}}\right)
        $$
        \item 由此我们成功完成了对期望的拆分，可以来考虑求出每个位置的系数是多少
    \end{itemize}
\end{frame}

\begin{frame}{Solution}
    求系数
    \begin{itemize}
        \item 因为是在求系数，我们完全不关心 $a_i$，所以自然我们只关心左侧和右侧有多少个数
        \item 这给了我们一个 $n^2$ DP 的尝试，设 $f_{i,j}$ 表示左侧有 $i$ 个数，右侧有 $j$ 个数，此位置最终系数的期望是多少。那么 $k_i=f_{i-1,n-i}$
        \item 转移只需要尝试在最左侧或最右侧加一个数$n^2$ DP 的尝试，设 $f_{i,j}$ 表示左侧有 $i$ 个数，右侧有 $j$ 个数，此位置最终系数的期望是多少。那么 $k_i=f_{i-1,n-i}$

        $$
            f_{i,j}=\frac{i-1}{i+j}\times f_{i-1,j}+\frac{j-1}{i+j}\times f_{i,j-1}+\frac{1}{i+j}\times \frac{f_{i,j-1}}{2}+\frac{1}{i+j}\times \frac{f_{i-1,j}}{2}
        $$
    \end{itemize}
\end{frame}

\begin{frame}{Solution}
    求系数
    \begin{itemize}
        \item 状态为 $n^2$ 没有前途，尝试优化
        \item 我们发现同时对两侧做 DP 非常蠢，他们实际上是独立的
        \item 因此只对一侧做，设 $g_i$ 表示右侧有 $i$ 个数字，最后只剩下一个时，第一个数的系数期望
        \item 转移依旧考虑加一个数只需要考虑最后一个空隙被合并的顺序，$g_i=\frac{1}{i}\times\frac{g_{i-1}}{2}+\frac{i-1}{i}\times g_{i-1}$
        \item 则 $k_i=g_{i-1}\times g_{n-i}$
        \item 复杂度瓶颈是求逆元，可以线性求逆元，也可以直接快速幂，均可快速通过
    \end{itemize}
\end{frame}

\section{Taki 的融合}
\subsection{Taki 的融合}
\begin{frame}{题面}
    若干序列，每次可以任何一个序列的开头取一个数加入操作序列 $opt$，要求最小化

    $$
        \sum_i \sum_{i\le j} opt_j
    $$

    多次修改，每次修改形如在某个序列开头或结尾添加元素

    $\sum len,k,q\le 3\times 10^5$
\end{frame}

\begin{frame}{Solution}
    \begin{itemize}
        \item 先不考虑带修的情况，对式子变形，考虑每个数的贡献，会得到
        $$
            \sum_{i=1}^{N}(N+1-i)opt_i
        $$
        \item 考虑序列内部，若 $x<y$，则 $x$ 先取显然更优，即 $2x+y<2y+x$
        \item 但有可能 $x$ 排在 $y$ 后面，而规则限制我们必须先取前列的数。因此若 $a_{i}>a_{i+1}$，我们一定会在取完 $a_i$ 之后立刻取 $a_{i+1}$，否则我们交换操作顺序一定不劣
    \end{itemize}
\end{frame}

\begin{frame}{Solution}
    \begin{itemize}
        \item 因此我们进一步采取交换论证，现在有两个 “连续取” 段 $(a_1,a_2,\cdots, a_n)$ 和 $(b_1,b_2,\cdots, b_m)$。我们有两种选择：
        \begin{itemize}
            \item 先 $a$ 后 $b$：$X=\sum_{i=1}^n a_i(n+m+1-i)+\sum_{i=1}^mb_i(m+1-i)$
            \item 先 $b$ 后 $a$：$Y=\sum_{i=1}^m b_i(n+m+1-i)+\sum_{i=1}^na_i(n+1-i)$
        \end{itemize} 
        \item 作差 $X-Y=m\sum a_i-n\sum b_i$，再令 $X-Y<0$ 得到 $\frac{\sum_i a_i}{n}<\frac{\sum b_i}{m}$
        \item 这表明先 $a$ 后 $b$ 当且仅当 $a$ 的平均数更小
        \item 再来思考合并连续段需要维护什么信息：$X=m\sum_{i=1}^n a_i+\sum_{i=1}^n a_i(n+1-i)+\sum_{i=1}^mb_i(m+1-i)$
        \item 可以发现，只需要维护各自的答案、总和和元素个数即可
        \item 由此我们确定了我们的策略：先不断合并序列内部，直到序列内变为若干连续段，满足平均数不降，再对所有连续段排序，从小到大合并
    \end{itemize}
\end{frame}

\begin{frame}{Solution}
    \begin{itemize}
        \item 现在考虑带修，显然在序列前端和后端加数字只会促进段的合并。例如，有可能在开头加了一个较大的元素，这使得一些前缀发生了合并
        \item 既然只会发生合并，那么连续段的总个数是 $O(n+q)$ 的
        \item 可以直接用平衡树实时维护有哪些连续段。
        \item 更好的做法是离线出所有出现过的连续段，记录每个连续段出现与消失的时间。最后把出现过的所有连续段排序，再挂到线段树叶子上，在扫描线的过程中用 push\_up 维护合并，根节点的值即为答案
        \item 复杂度为 $O((n+q)\log(n+q))$，可以使用 std::deque 或者 std::list 维护，具体实现可以参考 std。
    \end{itemize}
\end{frame}






% ===== 结束页 =====

\begin{frame}
    \begin{center}
        {\Huge\calligra Thanks!}
    \end{center}
\end{frame}

\end{document}
